% =====================================================================
\section{Introduction}
\label{sec:intro}
% =====================================================================

Every day at noon, bids from several hundred generators, retailers, industrial consumers and traders are collected for an auction that sets the next day's electricity prices across the Nordic and Baltic region. The resulting area prices, which Norwegian households on spot-price contracts pay, are calculated in the European market coupling and published by Nord Pool. In 2025, Nord Pool traded 1,231~TWh across its markets \citep{NPNewsletter2026}. In 2024, 400 businesses from 20 countries traded on its platforms, and the average Nordic system price was €36.06/MWh \citep{NP2024}.

This case study analyses Nord Pool not as a single firm but as the platform of a cross-sector alliance. Since January 2020, Nord Pool has been owned 66\% by Euronext, a listed pan-European operator of financial market infrastructure, and 34\% by TSO Holding, a joint company owned by state-owned electricity transmission system operators (TSOs) \citep{Euronext2020,NPAbout}. TSO Holding is currently owned by EPSO-G of Lithuania (39.6\%), Statnett of Norway (32.2\%) and Svenska kraftnät of Sweden (28.2\%) \citep{NPAbout}. The alliance therefore joins two sectors with very different objectives: a profit-seeking exchange group and public-interest grid operators that are responsible for security of supply.

The alliance meets the three case criteria. First, the partners are independent organisations from different sectors that share ownership and governance of the platform \citep{NPAbout,NPBoard}. Second, Nord Pool provides a data-intensive, platform-based service for critical infrastructure: price formation, clearing and market data for electricity. Third, its ownership, fees, governance and market activities are well documented in public sources.

The alliance has also changed shape over time, which makes it an instructive case. Nord Pool was created in 1996 as the world's first multinational power exchange \citep{NPAbout}, building on Norway's market liberalisation in 1991 \citep{AmundsenBergman2006}. It remained fully owned by the TSOs until Euronext acquired control \citep{Euronext2020}, and in 2022 the Finnish and Danish TSOs sold their stakes in TSO Holding \citep{EPSOG2022}. Each change shifted who controls the platform, who contributes to it and who carries its risks.

We argue that the alliance has separated the actors who \emph{contribute} to the Nordic power price from the actors who \emph{control} it. The report develops this argument through five analytical tasks. Section~\ref{sec:t1} maps the business ecosystem. Section~\ref{sec:t2} analyses network effects and path dependency. Section~\ref{sec:t3} treats Nord Pool as a multisided platform and examines its pricing. Section~\ref{sec:t4} analyses cybersecurity as a good provided jointly by the partners, and Section~\ref{sec:t5} presents the business model and stakeholder relationships. Section~\ref{sec:conclusion} concludes by identifying the alliance's central strategic tension. Appendices~A and~B document our use of AI and are followed by a closing group reflection.
